\subsection{Definition Usability nach ISO 9241-11}
\label{ch:SdF:UI:sec:Usability}
Der Begriff der Usability (dt. Gebrauchstauglichkeit) beschreibt die Qualität eines Systems oder einer Oberfläche in Bezug auf dessen Einfachheit und Nutzbarkeit. Diese wird gemessen in den Dimensionen der Effektivität, Effizienz und der Zufriedenheit der Nutzer.

Die Effektivität beschreibt die funktionelle Eignung eines Systems. Das heißt, dass der Nutzer seine Aufgaben vollständig und korrekt mithilfe dieses Systems durchführen kann. Die Effizienz misst sich in Form von Zeit, um eine Aufgabe vollständig und richtig zu erfüllen, und der Rate, wie viele Fehler auf dem Weg dorthin gemacht werden. Die Zufriedenheit oder auch Akzeptanz (en. Satisfaction) \cite{Vollrath:2015:Taschenbuch} der Nutzer ist nicht in metrischen Kenngrößen zu erfassen. Es kann ein Zufriedenheitsspektrum auf Grundlage von zum Beispiel Fragebögen (s. \autoref{ch:SdF:UI:sec:ISO-Fragebogen}) oder der Verhaltensanalyse während eines Testszenarios festgestellt werden. [Zusammenfassung nach \cite{DIN_EN_ISO_9241-110}]

% ******************************************************************
\subsection{Definition User-Experience und Abgrenzung von Usability}
\label{ch:SdF:UI:sec:Usability_vs_User_Experience}
Neben der Usability steht die User-Experience (UX, dt. Nutzererfahrung oder -erlebnis). Dessen Inhalte leicht verwechselt werden. Die User-Experience wird in Part 210 der ISO 9241 definiert als \glqq Wahrnehmungen und Reaktionen einer Person, die aus der tatsächlichen und/oder der erwarteten Benutzung eines Produktes, eines Systems oder einer Dienstleitung resultieren.\grqq{} \cite[S. 7]{DIN_EN_ISO_9241-210:2011}

Während bei der Usability die schnelle, fehlerfreie und angenehme Bedienung eines Systems im Vordergrund stehen, wird die User-Experience zusätzlich mit Spaß, Vergnügen und positivem Erleben verbunden \cite{Vollrath:2015:Taschenbuch}.

Da die Gebrauchstauglichkeit (UE) ein Part des Nutzererlebnisses (UX) ist ($UX \varsubsetneq UE$) \cite{usability_vs_experience:2018} wird diese nicht in \autoref{fig:iso-uebersicht} aufgelistet. Der Part 112 mit den Grundsätzen zur Informationsdarstellung ist vor allem auf Benutzeroberflächen oder Dash\-boards anwendbar. Da sich dieser Part auf UX fokussiert, finden dessen Prinzipien hier aus Umfangsgründen keine Anwendung.

